% !PRINTED 400
\weberpage{400}
%<<CONT
den positiven Werthen von $t$, $u$; eine zweite ist positiv und
kleiner als 1 und die beiden anderen sind negativ. Ist $T$, $U$
die kleinste positive Lösung der Pell'schen Gleichung, so ist
\[
\Theta = \frac{T + U\sqrt{D}}{2}
\]
die kleinste unecht gebrochene zu $\sqrt{D}$ gehörige positive Einheit,
und wir können nachweisen, dass in der Form $\pm \Theta^n$, worin $n$
eine ganze Zahl bedeutet, alle zu $\sqrt{D}$ gehörigen Einheiten enthalten
sind. Es genügt dazu, zu zeigen, dass der Ausdruck $\Theta^n$
für ein positives $n$ alle Einheiten, die grösser als 1 sind, liefert.
Dies ist aber sehr einfach; denn ist $\Theta_1$ irgend eine von diesen
Einheiten, so wird sie zwischen zwei auf einander folgenden
Potenzen von $\Theta$ liegen, da diese Potenzen mit dem Exponenten
ins Unendliche wachsen; also
\[
\Theta^n \leqq \Theta_1 < \Theta^{n+1},
\]
mithin
\[
1 \leqq \Theta_1 \Theta^{-n} < \Theta.
\]
Da aber $\Theta_1 \Theta^{-n}$ gleichfalls eine Einheit ist, so ist, wenn nicht
das Gleichheitszeichen gilt, $\Theta_1 \Theta^{-n}$ grösser als 1 und kleiner als
$\Theta$, was gegen die Voraussetzung ist, dass $\Theta$ die kleinste unecht
gebrochene Einheit sei; es muss also
\[
\Theta_1 = \Theta^n
\]
sein, was zu beweisen war.

\gesperrt{Anmerkung, die \gesperrt{Gauss}'sche Theorie der quadratischen
Formen betreffend}\footnote{Vergl. Disq. ar. art. 183 ff., \gesperrt{Dirichlet}-\gesperrt{Dedekind}, Vorlesungen
über Zahlentheorie. Vierte Auflage. \S.\,72 ff.}.

Wir fügen, um die Verbindung unserer Betrachtungen mit
der sonst bekannten \gesperrt{Gauss}'schen Theorie der quadratischen
Formen herzustellen, einige Bemerkungen bei, die in unserem
Zusammenhange nicht gerade erforderlich sind, und die daher
auch übergangen werden können.

Zwei quadratische Formen,
\[
(A, B, C), \quad (-A, B, -C)
\]
nach der Bezeichnung von \gesperrt{Gauss}, sind im \gesperrt{Gauss}'schen Sinne
dann und nur dann mit einander äquivalent, wenn die Periode
%CONT>>
